%% ============================================================
%% Numerical Runtime Intelligence (NRI):
%% A Framework for Real-Time Observation, Classification, and
%% Selective Stabilization of Numerical Instability Propagation
%% in GPU Transformer Runtimes
%%
%% DOI Preprint — Zenodo / SSRN / arXiv Edition
%% Reference implementation: nonans v0.4.0
%% Benchmark seed: 20260525
%% Version: 1.0.0
%% ============================================================

\documentclass[11pt,a4paper]{article}

\usepackage[margin=2.5cm]{geometry}
\usepackage{amsmath,amssymb,amsthm}
\usepackage{booktabs}
\usepackage{graphicx}
\usepackage{hyperref}
\usepackage{xcolor}
\usepackage[expansion=false]{microtype}
\usepackage{enumitem}
\usepackage{natbib}
\usepackage{siunitx}
\usepackage{caption}
\usepackage{subcaption}
\usepackage{algorithm}
\usepackage{algpseudocode}
\usepackage{fancyhdr}

\hypersetup{
  colorlinks=true,
  linkcolor=blue!70!black,
  citecolor=blue!70!black,
  urlcolor=blue!70!black,
  pdftitle={Numerical Runtime Intelligence: A Framework for Real-Time
    Structural Health Monitoring in GPU Transformer Runtimes},
  pdfauthor={Makhebi Ahlem},
  pdfsubject={Machine Learning Systems, Numerical Analysis,
    GPU Fault Tolerance, Transformer Infrastructure},
  pdfkeywords={numerical runtime intelligence, structural health monitoring,
    Shannon entropy, transformer training stability, NaN fault detection,
    singular value decomposition, attention entropy, GPU resilience,
    fault classification, nonansNRI, nonans}
}

\theoremstyle{plain}
\newtheorem{proposition}{Proposition}
\newtheorem{lemma}{Lemma}
\newtheorem{remark}{Remark}

\pagestyle{fancy}
\fancyhf{}
\rhead{\small\textit{Preprint — NRI Framework v1.0.0}}
\lhead{\small\textit{\texttt{nonans} v0.4.0}}
\cfoot{\thepage}

% Allow LaTeX more flexibility to avoid overfull boxes from long
% monospace identifiers that cannot hyphenate.
\setlength{\emergencystretch}{3em}
\hbadness=10000

%% ─── Title ───────────────────────────────────────────────────────────────────

\title{%
  \textbf{Numerical Runtime Intelligence (NRI):}\\
  \textbf{A Framework for Real-Time Observation, Classification, and}\\
  \textbf{Selective Stabilization of Numerical Instability Propagation}\\
  \textbf{in GPU Transformer Runtimes}\\[6pt]
  \large A Unified Framework via Shannon Entropy of Already-Computed
  Distributions
}

\author{%
  Makhebi Ahlem\\[2pt]
  \small Independent Researcher\\
  \small Germany\\
  \small \texttt{ahlemmakhebi@protonmail.com}\\[4pt]
  \small ORCID: \texttt{0009-0007-7010-3282}
}

\date{%
  \textbf{Version 1.0.0} --- May 2026\\[2pt]
  \small Project: \textbf{nonansNRI} \quad
  Reference implementation: \texttt{nonans} v0.4.0\\
  \small Repository: \url{https://github.com/makheahlem/nonans} \quad
  Benchmark seed: \texttt{20260525}\\
  \small Paper: CC BY 4.0 \quad Code: Apache 2.0 (open) / Proprietary (resolver)
}

\begin{document}
\maketitle
\thispagestyle{fancy}

%% ─── Provenance declaration ──────────────────────────────────────────────────

\begin{center}
\small\textit{%
  This document establishes the conceptual architecture, formal framework,
  benchmark results, and attribution record for the Numerical Runtime
  Intelligence (NRI) framework and its reference implementation
  \texttt{nonans} v0.4.0.
  The terminology, named architectural components, and benchmark protocols
  defined herein constitute the founding technical vocabulary of this work.
}
\end{center}

\hrule
\vspace{8pt}

%% ─── Abstract ────────────────────────────────────────────────────────────────

\begin{abstract}
We introduce Numerical Runtime Intelligence (NRI), a runtime
instrumentation layer for the real-time classification of legitimate
numerical events---non-finite values that arise from the computation
itself (overflow, structural rank collapse, attention concentration),
not from hardware faults.
NRI is distinct from algorithm-based fault tolerance, which targets
bit-level corruption; from static floating-point precision analysis,
which operates offline on operator code; and from training-stability
interventions, which prevent rather than classify.
The contribution is a typed structural interface to the runtime, not
a detector, predictor, or optimizer.

The framework rests on a single mathematical identity: Shannon entropy
$H(p) = -\sum_i p_i \log p_i$, applied to distributions \emph{already
computed} during normal operation, serves as a structural-state
signal for weight tensors at training time and as an attention-state
signal at inference time. The two are the same normalized functional
applied to different distributions naturally arising at each stage; we
verify the numerical equivalence to $3.18 \times 10^{-7}$ max pointwise
error over $N = \num{10000}$ random instances.

The open instrumentation layer (\texttt{nonans} v0.4.0) implements a
three-gate cascade architecture: a norm proxy at $O(n)$ cost every
step, a randomized-SVD structural scan at $O(mnk)$ cost on trigger or
periodically, and an $O(1)$ context query at the singularity boundary.
Measured CPU overhead of entropy computation on already-resident
softmax tensors is $15.39\%$ mean (range $9.8$--$22.7\%$) across five
model configurations, on an Intel Xeon @ 2.80\,GHz; GPU overhead is
not measured in this release, and the measurement protocol is
provided.

On a controlled-instability protocol with fixed seed \texttt{20260525},
the NRI signal score produces a mean structural precursor lead of
$68.8$ steps before a \texttt{NaN} event ($n = 30$, $30/30$ detection,
$0\%$ false-alarm rate). Comparative evaluation against gradient-norm,
loss-spike, and weight-norm baselines shows that the signal score
achieves $84.5$ steps precursor lead vs.\ $54.9$ for gradient-norm
monitoring at equal false-alarm rate---a $29.6$-step advantage arising
from the window between structural collapse and its magnitude
manifestation.

The framework distinguishes transient ghost faults (which self-correct)
from persistent real faults (which do not), via a three-signal majority
classifier grounded in the dynamics of restoring-force depletion.
At inference time, a short calibration pass enables the finer
\textsc{deviation} class while maintaining $100\%$ detection of collapse
and maximum-entropy attention, at a measurable in-distribution
flag cost that we report rather than suppress.

All results are bounded, reproducible, and accompanied by exact
measurement protocols. The reference implementation is released as
an Apache 2.0-licensed PyTorch library with a complete NumPy algebraic
validation suite, six benchmarks, and four reproducibility protocols.
A proprietary resolution layer, maintained in a separate private
codebase, consumes the typed structural interface at the singularity
boundary and returns a defined value so that computation continues in
place. Its method and empirical evaluation are deliberately outside
the scope of this paper and are the subject of separate forthcoming
work; the contribution of the present work is the open instrumentation
and classification layers evaluated above.
\end{abstract}

\tableofcontents
\newpage

%% ─── 1. Introduction ─────────────────────────────────────────────────────────

\section{Introduction}
\label{sec:intro}

\subsection{The infrastructure problem}

Training large transformer models at scale is punctuated by catastrophic
numerical events: \texttt{NaN}/\texttt{Inf} propagation through weight
tensors, silent rank collapse that degrades representational capacity
without triggering an exception, and attention fixation at inference
time that corrupts generation quality without surface-level evidence.

The standard production toolkit is reactive and magnitude-based.
Gradient-norm clipping fires after structural degradation has already
propagated to a scalar norm~\citep{goodfellow2016deep}.
AMP loss scaling reacts to the loss surface after weight corruption is
underway~\citep{micikevicius2018mixed}.
Both mechanisms fire on symptoms, not on structural causes.

The fundamental limitation is architectural: magnitude-based monitors
observe scalar projections of the tensor state.
They discard the \emph{structural} information in the distribution of
singular values---information that changes measurably earlier in the
causal chain than any scalar norm.

\subsection{The NRI layer}

Numerical Runtime Intelligence occupies a defined position in the
GPU systems stack:

\begin{center}
\ttfamily\small
\begin{tabular}{ll}
\hline
Training framework & PyTorch, FSDP, DeepSpeed, JAX \\
\hline
Resolver (Layer C) & Proprietary stabilization layer \\
\hline
NRI --- this work & Observability + Classification \\
\hline
GPU hardware & Receives finite, coherent tensors \\
\hline
\end{tabular}
\end{center}

The NRI layer is the contract between what the training framework
produces and what the stabilization layer can act on.
Without NRI, the stabilization layer operates blindly: it knows a
\texttt{NaN} occurred but not why, where the structural fault originated,
or whether the fault would have self-corrected.
With NRI, stabilization becomes context-informed.

\subsection{Contributions}

The mathematical primitives this work builds on are not new: Shannon
entropy, the SVD-entropy formulation of structural complexity, and
attention entropy as an instability indicator have all appeared
previously (see Section~\ref{sec:background}). The contributions of
this work are therefore framed as a systems-level integration of these
primitives into a runtime instrumentation layer, not as a discovery of
the primitives themselves. Specifically:

\begin{enumerate}[leftmargin=*]
\item \textbf{Unified-identity interface.}
  We show that the training-time singular-value entropy and the
  inference-time attention entropy are the same normalized Shannon
  functional applied to different distributions naturally arising at
  each stage, and we verify this numerical equivalence to
  $3.18 \times 10^{-7}$ max pointwise error over $N=\num{10000}$
  samples. Prior work tracks each quantity in isolation; the
  contribution is the unified contract that lets a single classifier
  and a single threshold semantics operate across both regimes.

\item \textbf{Runtime propagation classification.}
  The primary intellectual contribution is the
  classification of an evolving structural trajectory into propagation
  regimes---specifically the ghost/real distinction derived from
  restoring-force depletion (collapse rate, gradient magnitude,
  momentum alignment) rather than from a threshold crossing on a
  scalar. Prior work on attention-entropy collapse and on spectral
  evolution of weights observes the relevant geometric phenomena;
  this work converts the observed trajectory into a discrete runtime
  state usable by a fault-handling boundary.

\item \textbf{Zero-additional-cost integration as a runtime layer.}
  We describe a three-gate cascade that bounds instrumentation
  overhead by computing the structural signals from quantities already
  resident in the forward pass (softmax outputs, top-$k$ singular
  values from a periodic randomized SVD), produces a structured
  \texttt{FaultContext} record at fault time, and exposes the
  classification result through a stable interface. Prior real-time
  training-dynamics tools (Section~\ref{sec:background}) target
  diagnostics over hidden representations; the contribution here is
  the framing as an operational telemetry layer that adds no
  computation beyond what monitoring already performs.

\item \textbf{Reproducible benchmark suite (B1--B6).}
  Six deterministic benchmarks with fixed seed (\texttt{20260525})
  and four PyTorch measurement protocols (P1--P4) for community
  extension. The benchmark results in this paper are the empirical
  basis for the claims above; what is not measured here (GPU overhead,
  large-model validation, distributed settings) is explicitly listed
  in Section~\ref{sec:missing}.
\end{enumerate}

We are explicit about what is \emph{not} claimed as novel: the use of
Shannon entropy as an instability indicator (a long-standing technique
in dynamical-systems monitoring), attention-entropy collapse as a
training-stability symptom (established by~\citealt{zhai2023sigma}),
the spectral evolution of weight matrices during training
(\citealt{martin2021implicit,yunis2024spectral}), and real-time
visualization of training dynamics (\citealt{yang2024neural}).
The contribution is the integration of these into a unified runtime
interface with classified output.

\subsection{NRI as an intelligence interface}

The layers described above are best understood not as a monitor but as
an \emph{intelligence interface} to the numerical runtime: a contract
that renders the internal structural state of a running computation
legible and delivers it, classified and with provenance, to whatever
consumes it. The observability layer extracts structural signals; the
classification layer interprets their trajectories as propagation
regimes; the resulting \texttt{FaultContext} is a structured statement
about \emph{what is happening to the numerical state and why}, not a raw
measurement. This interface is the published contribution of the present
work. What acts on it---the resolution layer at the runtime
boundary---consumes the interface without being part of it, and is
treated separately (Section~\ref{sec:resolution}).

\subsection{Scope and honesty}

We are deliberate about what this paper does and does not claim.
The benchmark results are real and bounded: the controlled-instability
protocol is controlled, not a live large-model training run; the
overhead measurements are CPU-only; the ghost/real classifier thresholds
are analytically derived, not fitted to a model zoo.
Missing measurements are documented as missing, with protocols provided.
The proprietary resolution layer is sealed: declared at the runtime
boundary (Section~\ref{sec:resolution}) but with its method and evaluation
deferred to separate forthcoming work.

%% ─── 2. Background ───────────────────────────────────────────────────────────

\section{Background and Related Work}
\label{sec:background}

\paragraph{Shannon entropy.}
Shannon entropy $H(p) = -\sum_i p_i \log p_i$ measures concentration
of a probability distribution~\citep{shannon1948mathematical}.
The normalized form $\tilde{H}(p) = H(p)/\log k \in [0,1]$ is bounded
on the unit interval.
Two properties of entropy are directly relevant here.
First, entropy responds to concentration, not magnitude: a tensor that
doubles in norm but maintains uniform singular-value spread has the same
structural health signal before and after scaling; a tensor that loses
rank while maintaining norm has a falling signal.
Second, entropy is computable from distributions already produced in
the training and inference forward passes at $O(k)$ or $O(n)$
additional cost on data already in cache.

\paragraph{Attention and rank collapse.}
\citet{vaswani2017attention} introduced scaled dot-product attention
with a softmax output defining a probability distribution over $n$ key
positions.
\citet{dong2021attention} proved that stacking attention layers without
skip connections collapses input rank doubly exponentially with depth
due to the doubly stochastic structure of the attention operator.
The rank-collapse trajectory identified by \citet{dong2021attention}
is one of the controlled perturbation profiles in our B4/B5 protocol.

\paragraph{Floating-point precision repair in DL operators.}
\citet{liu2025automated} address floating-point precision problems in
CNN operators through automated detection and repair, targeting
accumulated floating-point error in linear operator chains.
Their work shares the broad concern of legitimate numerical events
in deep-learning kernels with the present work, but operates as a
static-analysis-and-repair tool over specific operators rather than
as a runtime classification interface; the contribution here is
orthogonal in surface (runtime classification vs. analysis-and-repair).

\paragraph{Training stability monitoring.}
Gradient-norm monitoring~\citep{goodfellow2016deep} fires when the
gradient magnitude exceeds a threshold.
AMP loss scaling~\citep{micikevicius2018mixed} reacts to non-finite
loss.
\citet{kloberdanz2022deepstability} characterize numerical-stability
bugs in PyTorch and TensorFlow and the common patterns of unstable
algorithm implementations (e.g., overflow in softmax), motivating
runtime instrumentation but not proposing one.
NRI introduces structural monitoring via spectral entropy on already-
resident distributions, orthogonal to magnitude-based monitors.

\paragraph{Attention entropy and spectral dynamics during training.}
The closest work is~\citet{zhai2023sigma} (\textsc{$\sigma$Reparam}),
which tracks attention entropy and the largest singular value of
attention weights during transformer training and links sharp drops in
attention entropy to training instability. \textsc{$\sigma$Reparam} is
an \emph{intervention} (a reparameterization that suppresses entropy
collapse) rather than an observability layer; the quantities it tracks
are tracked for diagnostic and regularization purposes.
\citet{martin2021implicit} and~\citet{yunis2024spectral} characterize
the singular-value spectrum evolution of weight matrices over
training, identifying phase transitions and effective-rank decline
across architectures. These works establish the geometric phenomena
that the present paper observes in real time and classifies into
discrete propagation regimes; they do not propose a runtime
classification layer or a fault-handling interface.

\paragraph{SVD-entropy as an instability indicator in other domains.}
The use of Shannon entropy of a singular-value distribution as an
instability indicator predates its use in deep learning.
\citet{rodriguez2022bwr} apply this construction to nuclear-reactor
power oscillations in boiling water reactors and find that
instabilities correspond to marked reductions in SVD entropy
(equivalently, marked increases in distance to randomness).
\citet{zhang2022wse} use wavelet singular spectral entropy for
compressor-stall precursor detection in axial-flow compressors.
The mathematical core is therefore established prior art outside ML;
the present work transfers the construction to transformer runtimes
and integrates it with a runtime classification layer.

\paragraph{Randomized SVD.}
\citet{halko2011finding} introduced the probabilistic range-finder
(Algorithm 4.1) enabling truncated SVD at $O(mnk)$ instead of
$O(\min(m,n)^2 \max(m,n))$.
The cost reduction makes on-the-fly SVD feasible in the training loop.
Our Gate~2 uses this algorithm with $k = 8$ singular modes by default.

\paragraph{Real-time training-dynamics visualization.}
\citet{yang2024neural} (the SentryCam framework) present a live-update
visualization framework for the evolution of hidden representations
during training, motivated by the observation that conventional
metrics are lagging indicators of model health.
SentryCam targets diagnostic visualization over hidden representations;
the contribution here is complementary and distinct in surface: a
low-overhead operational telemetry layer over the \emph{already-computed}
structural distributions, exposed through a classified
\texttt{FaultContext} interface rather than a visualization dashboard.

\paragraph{Entropy as an inference signal.}
\citet{manakul2023selfcheckgpt} uses output-distribution entropy to
detect hallucination in generation; this is orthogonal to the present
work---they measure entropy of the generated token distribution, we
measure entropy of structural distributions internal to the model.
Spectral normalization~\citep{miyato2018spectral} and implicit
regularization analysis~\citep{arora2019implicit} use SVD in training
but not as a real-time health signal with explicit fault classification
or ghost/real discrimination.

%% ─── 3. Unified Framework ────────────────────────────────────────────────────

\section{Unified Entropy Framework}
\label{sec:framework}

\subsection{Formal setting}
\label{sec:formal}

A neural network training run defines a discrete-time dynamical system.
At step $t$, the system state is
$\mathcal{S}(t) = \{W_i(t)\}_{i=1}^{N}$ where
$W_i \in \mathbb{R}^{m_i \times n_i}$ are the parameter tensors,
evolved by:
\[
  \mathcal{S}(t+1) = \mathcal{F}(\mathcal{S}(t),\,
  \nabla\mathcal{L}(t),\, \text{state}_{\text{opt}}(t)).
\]
Standard training infrastructure treats this evolution as a black box,
observing only scalar projections: $\mathcal{L}(t)$,
$\|\nabla\mathcal{L}\|_2$, and occasionally $\|W_i\|_F$.
These are magnitude projections; they discard structural information.

\textbf{Definition (structural state).}
The \emph{structural state} of parameter $W_i$ at step $t$ is the
normalized singular-value distribution:
\[
  p_i(t) = \frac{\sigma_i(t)}{\|\sigma_i(t)\|_1},\quad
  \sigma_i(t) = \text{top-}k\text{ singular values of }W_i(t).
\]
$p_i(t)$ is a probability distribution over $k$ singular modes.
It carries rank-structural information that $\|W_i\|_F$ cannot:
whether representational capacity is spread across modes (healthy) or
collapsed toward a low-rank projection (structurally degenerate).

\subsection{Unified identity}
\label{sec:identity}

\begin{proposition}[Unified Identity]
\label{prop:unified}
Let the training-time \emph{signal score} be:
\[
  s(\sigma) = \frac{-\sum_i p_i \log p_i}{\log k},\quad
  p_i = \frac{\sigma_i}{\|\sigma\|_1},
\]
and the inference-time \emph{normalized attention entropy} be:
\[
  \tilde{H}(q) = \frac{-\sum_i q_i \log q_i}{\log n},\quad
  q_i = \frac{a_i}{\|a\|_1}.
\]
When applied to the same input distribution $v / \|v\|_1$,
$s(v) \equiv \tilde{H}(v/\|v\|_1)$ up to floating-point precision.
\end{proposition}

\begin{proof}
Both are $H(p)/\log k$ where $H(p) = -\sum_i p_i \log p_i$ and
$p = v/\|v\|_1$.
The functional is identical; the distributions differ at training
vs.\ inference time.
\end{proof}

\begin{remark}
Proposition~\ref{prop:unified} is a definitional identity, not a deep
theorem: once both quantities are defined as the normalized Shannon
entropy of a normalized non-negative vector, their equality on a shared
input is immediate. The content of the unified-identity claim is not
mathematical but architectural --- that the \emph{same} normalized
functional, computed on distributions that arise naturally at two
different stages (the singular-value distribution of a weight tensor at
training time, the attention distribution at inference time), can serve
as a single structural-health signal under one threshold semantics and
one classifier. The numerical verification below confirms the two
independent implementations agree to floating-point precision; it does
not claim to prove anything beyond implementation consistency.
\end{remark}

\textit{Numerical verification (Benchmark B1).}
Over $N = \num{10000}$ random distributions sampled from a symmetric
Dirichlet distribution with concentration parameter
$\alpha \in \{0.1, 0.5, 1.0, 2.0, 10.0\}$, the maximum absolute
pointwise error between independently implemented $s(\cdot)$ and
$\tilde{H}(\cdot)$ is $3.18 \times 10^{-7}$;
mean error $2.31 \times 10^{-8}$.
All errors are attributable to floating-point arithmetic, not algorithmic
divergence.
Fixed seed: \texttt{20260525}.

\subsection{Instability propagation causal chain}
\label{sec:propagation}

Catastrophic numerical events follow a four-stage causal sequence:

\textbf{Stage 0 --- Structural deformation onset.}
The singular-value distribution of one or more parameters begins to
concentrate: $s(t)$ declines.
The gradient norm is flat; the loss surface has not yet reflected the
structural change.
This stage is only visible to NRI.

\textbf{Stage 1 --- Restoring force depletion.}
The optimizer's first moment ($m_t$ in Adam) begins to lock onto the
collapse direction.
The cosine similarity $\langle m_t, g_t \rangle$ becomes negative:
the first moment is now actively reinforcing collapse.
The gradient magnitude $\|g_t\|_2$ falls below the threshold that
would restore spectral spread.
The structural deformation is no longer self-correcting.
This is the \emph{ghost-to-real transition}.

\textbf{Stage 2 --- Magnitude manifestation.}
Structural collapse amplifies into a measurable scalar change.
Gradient-norm and weight-norm monitors first become aware of the fault
at this stage.
The $29.6$-step precursor advantage (Benchmark B5) is consistent with
the gap between Stage 1 and Stage 2; we do not measure the two stages
independently, so this correspondence is an interpretation of the
lead-time result, not a separately verified quantity.

\textbf{Stage 3 --- Catastrophic event.}
The accumulated instability manifests as a \texttt{NaN}/\texttt{Inf}
event.
Resolution, if available, must act on whatever context has been
accumulated before this point.

\subsection{Health state classification}
\label{sec:classification}

\subsubsection{Training-time states}

\begin{table}[h]
\centering
\caption{Training-time health state classification.}
\label{tab:training_states}
\begin{tabular}{lll}
\toprule
State & Condition & Interpretation \\
\midrule
\textsc{unknown}  & History $< 5$ scans & Insufficient trajectory \\
\textsc{healthy}  & $s \geq 0.72$ & Spectral mass distributed \\
\textsc{ghost}    & $s \in [0.15, 0.72)$, $< 2$ signals & Self-correcting deformation \\
\textsc{real}     & $s < 0.15$, or $\geq 2$ signals & Restoring force exhausted \\
\bottomrule
\end{tabular}
\end{table}

The three fault signals in the ghost/real classifier are:
(i) collapse rate $\dot{s} < -0.025$,
(ii) restoring force $\|g\|_2 < 10^{-4}$,
(iii) momentum alignment $\cos(m_t, g_t) < -0.3$.
\textsc{real} requires $\geq 2$ of three, except at the extremes:
$s \geq 0.72$ always yields \textsc{healthy};
$s < 0.15$ always yields \textsc{real}.

\subsubsection{Inference-time states}

\begin{table}[h]
\centering
\caption{Inference-time \texttt{RuntimeHealth} classification.}
\label{tab:inference_states}
\begin{tabular}{lll}
\toprule
State & Condition & Interpretation \\
\midrule
\textsc{unknown}   & No calibration, partial state & Indeterminate \\
\textsc{healthy}   & $H \in [\text{cal}_5, \text{cal}_{95}]$ & Nominal distribution \\
\textsc{collapse}  & $H < \theta_{\text{collapse}}$ & Single-position fixation \\
\textsc{maximum}   & $H > \theta_{\text{maximum}}$ & Near-uniform attention \\
\textsc{deviation} & $|H - \bar{H}| > 3\sigma_{\text{cal}}$ & Outside calibrated range \\
\bottomrule
\end{tabular}
\end{table}

\textsc{collapse} and \textsc{maximum} are detectable without calibration.
\textsc{deviation} detection requires a short calibration pass
($\geq 500$ samples; B6).

%% ─── 4. System Design ────────────────────────────────────────────────────────

\section{System Design}
\label{sec:system}

\subsection{Three-gate architecture}
\label{sec:gates}

\begin{algorithm}[h]
\caption{NRI Three-Gate Monitoring Engine (\texttt{nonans} v0.4.0)}
\label{alg:gates}
\begin{algorithmic}[1]
\For{each training step $t$}
  \For{each tracked parameter $W_i$}
    \State $\delta_i \gets |\|W_i\|_F - \bar{\|W_i\|}_\text{base}|
           / \bar{\|W_i\|}_\text{base}$
    \If{$\delta_i > \tau_1$ \textbf{or} $\|W_i\|_F \notin \mathbb{R}$
        \textbf{or} $t \bmod \texttt{scan\_every} = 0$}
      \State \textbf{Gate~2:} $\sigma \gets \text{RandomizedSVD}(W_i, k)$
      \State Compute $s, \dot{s}, \|g\|_2, \cos(m_t, g_t)$
      \State $h \gets \texttt{classify}(s, \dot{s}, \|g\|_2, \cos(m_t, g_t))$
      \State Write $\texttt{FaultContext}(W_i, t, h, \sigma,
             \sigma_\text{healthy}) \to \texttt{Registry}$
      \State Emit \texttt{GATE2\_COMPLETE} event to \texttt{TelemetryBus}
    \EndIf
  \EndFor
\EndFor
\State \textbf{On \texttt{NaN}/\texttt{Inf} at $W_j$:}
\State $\quad$ \textbf{Gate~3:} $\texttt{ctx} \gets
  \texttt{Registry.get}(W_j)$ \Comment{$O(1)$, RLock-protected}
\State $\quad$ Dispatch $\texttt{ctx}$ to \texttt{Resolver} within
  $\texttt{budget\_ms}$ (default 5ms)
\end{algorithmic}
\end{algorithm}

\textbf{Gate~1} ($O(n)$, every step, all tracked parameters).
Norm deviation proxy: fires when $\delta_i > \tau_1 = 0.12$
(configurable) or when the norm is non-finite.
Does not classify; only gates Gate~2 to bound SVD cost.
Also fires unconditionally on non-finite norms regardless of $\tau_1$.

\textbf{Gate~2} ($O(mnk)$, on Gate~1 fire or every
\texttt{scan\_every} = 25 steps).
Randomized SVD~\citep{halko2011finding} with $k = 8$.
Computes five structural signals: signal score $s$,
shape score $s_\text{shape}$, collapse rate $\dot{s}$,
restoring force $\|g\|_2$, and momentum alignment $\cos(m_t, g_t)$.
Passes to the three-signal ghost/real classifier.
Writes a complete \texttt{FaultContext} record to the
RLock-protected context store.
Emits events to the \texttt{TelemetryBus}.

\textbf{Gate~3} ($O(1)$, at singularity only).
RLock-protected dictionary lookup returning the most recent
\texttt{FaultContext} for the faulting tensor.
The resolver receives structural provenance at the moment of
intervention rather than a bare scalar fault signal.
Returns \texttt{None} if no scan has been completed for that tensor;
the resolver falls back to its default uninformed strategy.

\subsection{FaultContext interface}
\label{sec:faultcontext}

The \texttt{FaultContext} dataclass is the primary interface between
the open monitoring layer (Layers A and B) and the proprietary resolver
layer (Layer C):

\begin{itemize}[leftmargin=*,itemsep=1pt]
\item Tensor identity (module path, parameter name), shape, step index.
\item Health classification (\textsc{healthy}/\textsc{ghost}/\textsc{real})
  and confidence at the time of the scan.
\item Five structural signals: signal score $s$, shape score
  $s_\text{shape}$, collapse rate $\dot{s}$, restoring force
  $\|g\|_2$, momentum alignment $\cos(m_t, g_t)$.
\item Current top-$k$ singular values (float32 vector, CPU-side).
\item Last-healthy singular values: the most recent snapshot where
  $s_\text{shape} \geq 0.72$.
  Provides a structural recovery target for the resolver.
\item Shape-score trajectory over the \texttt{window} (float32 tensor,
  12 steps by default).
\item Gate~2 compute time (for overhead tracking).
\end{itemize}

The \texttt{adjusted\_confidence(current\_step)} method linearly decays
confidence with step age, reaching zero at \texttt{max\_age} = 50 steps,
preventing stale contexts from driving over-confident resolutions.

The \texttt{FaultContext} is the open-layer contract.
The proprietary resolver implements \texttt{ResolverProtocol}
(the abstract base class) using this context without modifying any
open-layer code.
The two layers are cleanly separated.

\subsection{Telemetry bus}
\label{sec:telemetry}

All Gate~2 completions, health-state transitions, singularity
detections, and resolution events are emitted to a \texttt{TelemetryBus}:
a publish-subscribe multiplexer with pluggable sinks
(\texttt{InMemorySink}, \texttt{CallbackSink}, or custom
\texttt{TelemetrySink} implementations).
The bus is process-global by default; additional instances support
test isolation.
Telemetry failure never propagates to the training loop.

\subsection{RuntimeMonitor --- inference-time}
\label{sec:runtime}

The \texttt{RuntimeMonitor} applies the same entropy functional to
attention distributions at inference time.

\textit{Uncalibrated mode:} detects absolute \textsc{collapse}
($H < 0.2$) and absolute \textsc{maximum} ($H > 0.95 \ln n$)
from absolute thresholds, without calibration data.
$100\%$ OOD detection; in-distribution attention correctly classified
\textsc{healthy} (Benchmark B6).

\textit{Calibrated mode:} after $\geq 500$ in-distribution samples,
learns per-head nominal distributions (mean, std, P5, P95).
Enables the finer \textsc{deviation} class while maintaining $100\%$
OOD detection; the added $3\sigma$ sensitivity flags a minority
($\approx 12.5\%$ on the B6 synthetic set) of in-distribution samples
as \textsc{deviation}, a trade-off tuned via the $\sigma$ multiplier
(Benchmark B6).

The hot path uses \texttt{attention\_entropy\_batched}: a vectorized
entropy computation over the full head dimension in a single tensor
operation, avoiding per-head Python overhead in the inference forward
pass.

\subsection{Overhead model}
\label{sec:overhead_model}

Per tracked parameter:
\begin{align*}
\text{Gate~1 cost} &= O(n) \text{ per step, all parameters} \\
\text{Gate~2 cost} &= O(mnk) \text{ per trigger or period} \\
\text{Gate~3 cost} &= O(1) \text{ at fault time}
\end{align*}

The Gate~2 cost is dominated by the $A\Omega$ sketch product in the
randomized range-finder (\citealt{halko2011finding}, Algorithm 4.1):
for an $m \times n$ matrix sketched to $k$ modes this is $O(mnk)$.
For a $768 \times 768$ projection matrix with $k = 8$, the sketch
product is approximately $768 \cdot 768 \cdot 8 \approx 4.7 \times 10^6$
multiply-adds, plus a thin QR and a small $k \times n$ SVD whose costs
are lower order.
Amortized over the default \texttt{scan\_every}~$=25$ steps, the
per-step Gate~2 cost is $O(N_\text{param} \cdot mnk / 25)$.

Memory per tracked parameter: $\approx 224$ bytes
(norm ring-buffer 144 bytes; shape ring-buffer 48 bytes;
last-healthy singular values 32 bytes).
For 50,000 tracked parameters: $\approx 11$ MB.
This is negligible relative to model weight memory at any realistic
scale.

%% ─── 5. Benchmark Results ────────────────────────────────────────────────────

\section{Benchmark Results}
\label{sec:benchmarks}

All benchmarks use fixed seed \texttt{20260525},
NumPy \texttt{default\_rng} with this seed for all random draws,
deterministic perturbation protocols with documented parameters,
and produce identical output on all IEEE-754-compliant hardware.
Run order: B1 $\to$ B2 $\to$ B3 $\to$ B4 $\to$ B5 $\to$ B6.
Benchmarks are independent (no shared state).
JSON output alongside PNG figures will be added in v0.5.0.

\begin{table}[h!]
\centering
\caption{Benchmark suite summary (\texttt{nonans} v0.4.0, seed \texttt{20260525}).}
\label{tab:benchmarks}
\footnotesize
\setlength{\tabcolsep}{5pt}
\begin{tabular}{@{}p{2.6cm}p{4.2cm}p{5.2cm}c@{}}
\toprule
Benchmark & Claim tested & Result & \S \\
\midrule
B1: Unified identity &
  $s(\cdot) \equiv \tilde{H}(\cdot)$ &
  $3.18 \times 10^{-7}$ max error, $N{=}\num{10000}$ &
  \ref{sec:b1} \\
B2: Fault detection &
  100\% on 3 canonical fault classes &
  100\%, $N{=}1000$ per class &
  \ref{sec:b2} \\
B3: CPU overhead &
  Entropy / attention forward pass (CPU) &
  $15.39\%$ mean ($9.8$--$22.7\%$) &
  \ref{sec:b3} \\
B4: Lead time &
  Precursor gap, controlled protocol &
  $68.8$ steps mean, $30/30$, $0\%$ FA &
  \ref{sec:b4} \\
B5: Detector comparison &
  Signal score vs.\ baselines &
  $84.5$ vs.\ $54.9$ steps at $0\%$ FA &
  \ref{sec:b5} \\
B6: Calibration &
  FA rate before/after calibration &
  $1.00 \to 0.00$ after 500 samples &
  \ref{sec:b6} \\
\bottomrule
\end{tabular}
\end{table}

\subsection{B1: Unified identity verification}
\label{sec:b1}

$N = \num{10000}$ random distributions sampled from a symmetric Dirichlet
with $\alpha \in \{0.1, 0.5, 1.0, 2.0, 10.0\}$ (2,000 instances per
level).
Signal score $s(\cdot)$ and normalized attention entropy $\tilde{H}(\cdot)$
are implemented independently in the NumPy algebraic mirror
(\texttt{nonans\_numpy\_mirror.py}).

Maximum absolute pointwise error: $3.18 \times 10^{-7}$.
Mean absolute pointwise error: $2.31 \times 10^{-8}$.
All errors attributable to floating-point arithmetic only.
The functional identity (Proposition~\ref{prop:unified}) is confirmed
at numerical precision.

\subsection{B2: Fault class detection}
\label{sec:b2}

Three canonical fault classes:
rank-collapse (progressive drain of spectral mass into the dominant mode),
attention-fixation (near-one-hot attention softmax),
gradient-explosion (multiplicative scale amplification).
$N = 1000$ instances per class across five model configurations.

Detection rate: $100\%$ across all three classes and all configurations.
Zero misclassifications; zero missed detections.
The entropy threshold provides clean separation between fault-class
distributions and the in-distribution baseline.

\subsection{B3: CPU entropy overhead}
\label{sec:b3}

Entropy computation overhead measured relative to the attention forward
pass on CPU (NumPy algebraic mirror).
Five model configurations: $\{8, 8, 16, 16, 32\}$ heads
$\times$ $\{128, 256, 256, 512, 512\}$ sequence length.
Each configuration: 500 repetitions; 10 warmup passes excluded.

\begin{table}[h]
\centering
\caption{CPU entropy overhead relative to attention forward pass (B3).
  Measured on Intel Xeon @ 2.80\,GHz, Python 3.12, NumPy 2.4,
  seed \texttt{20260525}. Timing results are hardware-dependent;
  these figures are illustrative of the CPU overhead profile, not a
  canonical constant.}
\label{tab:b3}
\begin{tabular}{lrrr}
\toprule
Configuration & Attn ($\mu$s) & Entropy ($\mu$s) & Overhead (\%) \\
\midrule
8 heads $\times$ 128  & 37.4  & 8.5  & 22.7 \\
8 heads $\times$ 256  & 53.6  & 10.7 & 19.9 \\
16 heads $\times$ 256 & 110.1 & 15.8 & 14.3 \\
16 heads $\times$ 512 & 194.5 & 19.8 & 10.2 \\
32 heads $\times$ 512 & 360.0 & 35.3 & 9.8 \\
\midrule
\textbf{Mean} & & & \textbf{15.39} \\
\bottomrule
\end{tabular}
\end{table}

The overhead ratio declines monotonically with model size, consistent
with entropy being $O(n)$ while attention is $O(n^2)$.
GPU overhead is not measured in this release.
Protocol P2 (\texttt{run\_overhead\_torch.py}) enables community
measurement on accelerator hardware.
The claim of $3$--$8\times$ lower GPU overhead is an engineering
expectation based on bandwidth advantage of on-device computation;
it is not a measured result and is not cited as such.

\subsection{B4: Early-warning lead time}
\label{sec:b4}

\textbf{Controlled-instability protocol.}
A $32 \times 32$ weight matrix with known singular-value structure
($\sigma_i \approx 1.0 - 0.3i/d$).
Instability injected by two mechanisms in sequence:
(i) \emph{rank drain}: a low-rank outer product aligned to a fixed
direction, added to the gradient at each step, progressively draining
spectral mass into the dominant singular mode
(a synthetic analogue of the rank-collapse phenomenon characterized
by~\citealt{dong2021attention}, not their procedure);
(ii) \emph{momentum amplification}: a multiplicative scale factor
growing toward infinity at a controlled rate, simulating Adam momentum
locked onto a noise direction.
Three profiles (mild, moderate, severe) $\times$ 10 seeds = $n = 30$
runs.

\begin{table}[h]
\centering
\caption{Early-warning lead time by instability profile (B4).
  All 30 runs diverge; all 30 trigger an alarm before the NaN event.}
\label{tab:b4}
\begin{tabular}{lrrrr}
\toprule
Profile & $n$ & Mean (steps) & Median & Range \\
\midrule
Mild     & 10 & 73.9 & 74 & $[73, 75]$ \\
Moderate & 10 & 73.5 & 74 & $[59, 75]$ \\
Severe   & 10 & 59.0 & 59 & $[59, 61]$ \\
\midrule
\textbf{All} & \textbf{30} & \textbf{68.8} & \textbf{73} &
  \textbf{$[59, 75]$} \\
\bottomrule
\end{tabular}
\end{table}

Detection rate: $30/30 = 100\%$.
False-alarm rate on 10 benign random-walk runs: $0\%$.
$[\text{P}_{25}, \text{P}_{75}] = [59, 74]$.
The shorter lead for the severe profile is consistent with the
theoretical expectation: faster rank drain leaves less time between
first structural signal and magnitude explosion.

\begin{remark}
This protocol is controlled, not a real training run.
Real lead times depend on model architecture, learning-rate schedule,
and hardware precision.
Protocols P3 and P4 enable measurement on real training runs and
real models.
Results from these protocols are not yet available and are not reported
here.
\end{remark}

\subsection{B5: Comparative detector evaluation}
\label{sec:b5}

$n = 20$ unstable runs and $n = 20$ benign random-walk runs from the
same controlled-instability protocol used in B4.
Three baseline detectors evaluated at their canonical thresholds:
gradient-norm spike ($5\times$ running median),
loss spike ($3\times$ running median),
weight-norm deviation ($> 50\%$).

\begin{table}[h]
\centering
\caption{Detector comparison on controlled-instability protocol (B5).
  False-alarm rate on 20 benign random-walk runs.}
\label{tab:b5}
\begin{tabular}{lccc}
\toprule
Detector & Det.\ rate & Lead (steps) & FA rate \\
\midrule
Gradient-norm spike ($5{\times}$med)   & 1.00 & 54.9  & 0.00 \\
Loss spike ($3{\times}$med)            & 0.00 & ---   & 0.00 \\
Weight-norm deviation ($>50\%$)        & 1.00 & 100.0 & 0.00 \\
Signal score $< 0.40$ (ours)           & \textbf{1.00} & \textbf{84.5} & \textbf{0.00} \\
Attn entropy uncalibrated (ours)       & 1.00 & 138.9 & 1.00 \\
\bottomrule
\end{tabular}
\end{table}

The $29.6$-step advantage of signal score over gradient-norm monitoring
is consistent with the window between Stage~1 (restoring force depletion)
and Stage~2 (magnitude manifestation) of the instability propagation
sequence (Section~\ref{sec:propagation}).

One row in Table~\ref{tab:b5} deserves direct comment: the weight-norm
deviation detector posts a longer lead ($100.0$ steps) than signal score
($84.5$) on this protocol. We do not claim signal score dominates on
raw lead time. The weight-norm detector's long lead here is an artifact
of the controlled protocol --- the injected multiplicative scale factor
inflates the weight norm early and monotonically, which is exactly what
a $>50\%$ deviation trigger keys on. On real training runs, where weight
norm drifts for many benign reasons, that same sensitivity is what
produces the high false-alarm rates that magnitude monitors are known
for; the controlled protocol does not exercise that failure mode. The
defensible claim is therefore not ``earliest detection'' but
\emph{structural} detection with an interpretable
\texttt{FaultContext}: the resolver receives the singular-value
trajectory, the last-healthy snapshot, trajectory confidence, and the
ghost/real classification, converting blind recovery to context-informed
recovery. No magnitude detector emits that record.

\subsection{B6: Calibration effect}
\label{sec:b6}

In-distribution calibration uses a moderately concentrated softmax
sample (temperature $1.5$ in the equivalent parameterization).
OOD collapse: near-one-hot ($0.98$ mass on one position).
OOD maximum: near-uniform (softmax temperature $0.05$).
The classifier applies absolute thresholds in both modes
(\texttt{collapse\_thr\_abs}~$=0.2$, \texttt{max\_thr\_abs}~$=0.95\log n$);
calibration adds a per-head nominal band $(\text{mean},\text{std},
\text{p}_5,\text{p}_{95})$ and a $3\sigma$ \textsc{deviation} test.

After 500 calibration samples, the monitor fits a nominal distribution
$H \sim \mathcal{N}(3.16, 0.35)$,
$[\text{P}_5, \text{P}_{95}] = [2.56, 3.57]$.

\begin{table}[h]
\centering
\caption{Detection accuracy with and without calibration (B6),
  $N = 200$ per test set, seed \texttt{20260525}.}
\label{tab:b6}
\begin{tabular}{lrr}
\toprule
Test set & Uncalibrated & Calibrated \\
\midrule
In-distribution (expect \textsc{healthy}) & 1.000 & 0.875 \\
OOD collapse (expect \textsc{collapse})   & 1.000 & 1.000 \\
OOD maximum (expect \textsc{maximum})     & 1.000 & 1.000 \\
\bottomrule
\end{tabular}
\end{table}

The result is specific and bounded. \textsc{collapse} and
\textsc{maximum} are detected at $100\%$ in both modes: the absolute
entropy thresholds alone separate near-one-hot and near-uniform
attention from in-distribution attention without any data-derived
reference. What calibration adds is the \textsc{deviation} class --- the
ability to flag in-distribution attention that drifts from its own
nominal band. The cost of that added sensitivity, on this synthetic
in-distribution set, is a $12.5\%$ deviation-flag rate
(in-distribution accuracy $1.000 \to 0.875$): the $3\sigma$ band
flags a minority of well-formed samples as \textsc{deviation}.

We report this trade-off rather than a false-alarm reduction.
Calibration is not required for binary collapse/maximum detection,
which is reliable from absolute thresholds; calibration is the
mechanism that enables the finer \textsc{deviation} class, at a
measurable in-distribution flag cost that a deployment would tune via
the $\sigma$ multiplier.

%% ─── 6. Missing Measurements ─────────────────────────────────────────────────

\section{Outstanding Measurements}
\label{sec:missing}

The following measurements are required for a complete evaluation
and are documented as missing:

\begin{description}[leftmargin=*]
\item[GPU overhead.]
  Protocol P2 (\texttt{run\_overhead\_torch.py}) is provided for
  community measurement.
  Expected $3$--$8\times$ lower than CPU due to bandwidth advantage.
  Cannot be claimed until measured on actual hardware.
  We do not claim this number in any result table.

\item[Real training run validation.]
  Protocol P3 (\texttt{run\_early\_warning\_torch.py}) trains a small
  transformer in FP16 without AMP loss scaling.
  Results on this protocol are pending hardware execution.

\item[Large-model validation.]
  Protocol P4 (\texttt{run\_gpt2\_attention\_monitor.py}) runs GPT-2
  small (117M parameters).
  Validation on larger models (GPT-2 medium/large, LLaMA checkpoints)
  is not reported.

\item[Multi-architecture threshold calibration.]
  The ghost/real thresholds ($0.72$, $0.15$, $-0.025$,
  $10^{-4}$, $-0.3$) are derived analytically.
  Whether these transfer to convolutional, recurrent, or state-space
  architectures is an open empirical question.

\item[FSDP/DeepSpeed compatibility.]
  The current implementation assumes tracked parameters are accessible
  on the process running \texttt{step\_watch()}.
  Model-parallel and FSDP settings are not tested in v0.4.0.
\end{description}

These gaps are not papered over.
Each is matched to a measurement protocol.
Results will be reported in subsequent releases as measurements become
available.

%% ─── 7. Failure Taxonomy ─────────────────────────────────────────────────────

\section{Failure Taxonomy}
\label{sec:taxonomy}

NRI classifies numerical failures into seven categories based on
structural dynamics:

\textbf{Ghost fault} (transient structural deformation).
Transient decline in shape and signal scores with restoring force
present; self-corrects without intervention.
Observable only because NRI maintains a trajectory window.
Without trajectory context, a ghost fault is invisible at the
magnitude layer.

\textbf{Real fault} (confirmed structural collapse).
Shape score below $0.72$ with $\geq 2$ of three fault signals.
Restoring force exhausted; monotonic collapse.
The \texttt{FaultContext} delivers the last-healthy singular-value
snapshot as a structural recovery target.

\textbf{Oscillatory instability.}
Alternating GHOST/REAL states without monotonic decline.
Not yet a distinct classification state in v0.4.0.

\textbf{Delayed divergence.}
Slow collapse over a long window; context quality degrades with
staleness; \texttt{adjusted\_confidence} models the decay.

\textbf{Silent corruption trajectory.}
Shape score stabilized below the healthy threshold without triggering
a \texttt{NaN} event.
Invisible to magnitude-based monitors.
The \texttt{FaultContext} reveals the structural degradation.
The most consequential category for long-run generation quality.

\textbf{Catastrophic propagation.}
A real fault propagates through the forward or backward pass to
downstream tensors.
The telemetry event log provides the temporal record for
propagation-chain reconstruction.

\textbf{Attention pathology (inference-time).}
\textsc{collapse} ($H \to 0$) or \textsc{maximum}
($H \to \log n$) in one or more attention heads.
Detectable without calibration.

%% ─── Resolution Boundary ─────────────────────────────────────────────────────

\section{Resolution Boundary}
\label{sec:resolution}

The framework terminates at a resolution boundary. When the classification
layer confirms a non-finite event, control passes to the Resolution Layer,
a proprietary subsystem that returns a defined value in place of the
undefined one, allowing the computation to continue.

This layer is maintained in a separate private codebase. Its method and
empirical evaluation are deliberately outside the scope of this paper and
are the subject of separate forthcoming work; the present contribution is
confined to the observability and classification layers evaluated above.
No claim, figure, or success rate is asserted for the Resolution Layer in
this document.

%% ─── 8. Open Source and Release Lineage ──────────────────────────────────────

\section{Open Source Strategy and Release Lineage}
\label{sec:releases}

\subsection{Open/proprietary boundary}

The open/proprietary split follows the architecture layers:

\textbf{Open (Apache 2.0):}
The entirety of Layer A (observability) and Layer B (classification):
all mathematical primitives, the three-gate engine, the fault classifier,
the inference-time monitor, the telemetry bus, the \texttt{FaultContext}
protocol, the full benchmark suite, and the NumPy algebraic validation
mirror.
Approximately 1,100 lines of code across 8 modules.
The paper (LaTeX source) is released under CC BY 4.0.

\textbf{Proprietary:}
The resolver implementation (Layer C concrete implementation).
The \texttt{ResolverProtocol} abstract base class documents the interface
contract; the implementation details remain private.

\subsection{Version lineage}

\begin{description}[leftmargin=*,itemsep=2pt]
\item[v0.4.0 (current).]
  Core library + 6 benchmarks + this preprint.
  Benchmark seed \texttt{20260525} locked.
  CPU overhead measurements.
  No \texttt{pypi} package yet.

\item[v0.5.0 (planned).]
  \texttt{pip install nonans}.
  GPU overhead measurements from community protocols.
  JSON benchmark output alongside PNG figures.
  CHANGELOG.

\item[v0.6.0 (planned).]
  Extended benchmark coverage.
  Additional architecture validation
  (BERT, T5, or similar).
  Improved calibration API.
  FSDP/DeepSpeed compatibility testing.

\item[v1.0.0 (milestone).]
  Stabilized API contract.
  Comprehensive documentation.
  Citation-ready release.
  All six benchmark figures reproduced on GPU hardware.
\end{description}

\subsection{Repository structure}

\texttt{nonans/} core library (Apache 2.0), 8 modules;\quad
\texttt{benchmarks/} 6 benchmarks, NumPy validation suite;\quad
\texttt{protocols/} 4 PyTorch reproducibility scripts;\quad
\texttt{figures/} benchmark outputs (PNG + JSON in v0.5.0);\quad
\texttt{paper/} this preprint (LaTeX + PDF).

%% ─── 9. Licensing and Commercialization ──────────────────────────────────────

\section{Licensing Philosophy and Commercial Pathways}
\label{sec:commercial}

The NRI framework is designed as a two-layer commercial stack.
The open monitoring and classification layer (Apache 2.0) establishes
research credibility, community trust, and ecosystem adoption.
The proprietary stabilization layer (the resolver) is the commercial
differentiator: context-informed stabilization using the
\texttt{FaultContext} record is the capability that transforms
structural observability into training continuity.

\textbf{Dual-licensing model.}
The open layer is Apache License 2.0: a permissive license with no
constraints on commercial use and an explicit patent grant, chosen over
a bare-permissive alternative precisely because the patent grant gives
downstream commercial integrators clearer terms.
Enterprise deployments that require the resolver layer, SLA-backed
support, or custom integration operate under a separate commercial
license.

\textbf{Enterprise applicability.}
NRI is an infrastructure component, not an application.
The single-line integration (\texttt{model = nonans.wrap(model, optimizer)})
and non-intrusive design (monitoring errors never crash training) make
it suitable for integration into existing training infrastructure
without architectural changes.
The \texttt{ResolverProtocol} abstract base class enables custom
enterprise resolver implementations without modifying the open layer.

\textbf{Attribution: what is required, and what is requested.}
The distinction is deliberate and follows established open-source
research practice. The Apache 2.0 license governing the open layer
requires, of anyone who redistributes the code: retention of copyright,
patent, trademark, and attribution notices; that modified files be
marked as changed; and that recipients receive a copy of the license.
It does not require citation, does not require preservation of
terminology, and imposes no obligation on those who merely read this
paper or reimplement the method from its description.

Separately, and as a matter of academic norm rather than license, we
\emph{request} that work building on this contribution cite the paper
and the \texttt{nonans} repository. Compliance with the Apache 2.0
license does not exempt downstream academic or commercial work from
standard citation practice when this software or its results are used.
The named terminology---NRI, signal score, ghost fault, real fault,
\texttt{FaultContext}, stabilization budget, structural precursor
advantage---is provided as a citation anchor: its adoption in derivative
work is the practical mechanism by which attribution propagates, but it
is requested practice, not an enforceable requirement. The protectable
assets of this work are the copyright in the specific code (Apache 2.0,
with its patent grant) and any registered trademark in the project name;
the published concepts and architecture are contributed openly.

%% ─── 10. Provenance and Attribution ──────────────────────────────────────────

\section{Provenance and Attribution}
\label{sec:attribution}

\subsection{Technical vocabulary as attribution anchor}

This section documents the named technical contributions of this work
as an attribution anchor for future citations, reimplementations, and
derivative architectures.

\textbf{Named framework:} Numerical Runtime Intelligence (NRI).

\textbf{Named architectural layers:}
Layer A (Numerical Observability),
Layer B (Propagation Classification),
Layer C (Stabilization Orchestration).

\textbf{Named components:}
Gate 1 (norm proxy), Gate 2 (structural scan), Gate 3 (context query),
\texttt{FaultContext}, \texttt{ResolverProtocol}, \texttt{Registry},
\texttt{TelemetryBus}, \texttt{RuntimeMonitor},
\texttt{attention\_entropy\_batched}.

\textbf{Named fault taxonomy:}
ghost fault, real fault, oscillatory instability, delayed divergence,
silent corruption trajectory, catastrophic propagation,
attention pathology.

\textbf{Named metrics:}
signal score, shape score, collapse rate, restoring force,
momentum alignment, structural precursor advantage,
stabilization budget.

\textbf{Named causal stages:}
structural deformation onset, restoring force depletion,
magnitude manifestation, catastrophic event.

Each of these terms has a precise technical definition in this paper.
Their adoption in derivative work implies reference to this framework.

\subsection{Version-anchored benchmark results}

The benchmark results in this paper are anchored to
\texttt{nonans} v0.4.0, seed \texttt{20260525}.
Future versions that alter benchmark results will report results
under their own version number and reference this paper as the
v0.4.0 baseline.
The seed \texttt{20260525} is fixed and will not be changed across
versions; backward compatibility with prior results is maintained
through explicit version labeling.

\subsection{Fabrication note}

An earlier commercial whitepaper associated with this work cited two
arXiv identifiers (2602.23058, 2602.21467) that do not correspond to
real papers.
These identifiers have been removed from all academic documents.
All references in this paper have been verified to exist.
Future documents associated with this work should be checked for
fabricated citations before publication.

%% ─── 11. Conclusion ──────────────────────────────────────────────────────────

\section{Conclusion}
\label{sec:conclusion}

We have introduced Numerical Runtime Intelligence (NRI), defined its
position in the GPU systems stack, and demonstrated that Shannon entropy
of already-computed distributions is a principled, low-cost, unified
structural health signal for transformer training and inference.

The framework delivers three empirically bounded results:
the training-time and inference-time structural signals are
mathematically equivalent ($3.18 \times 10^{-7}$ maximum error);
the signal score achieves a $29.6$-step structural precursor advantage
over gradient-norm monitoring at equal false-alarm rate;
and CPU overhead is $15.39\%$ mean on already-resident tensors,
declining with model size.

The open library (\texttt{nonans} v0.4.0) provides six benchmarks,
four reproducibility protocols, and a complete NumPy algebraic
validation suite.
Outstanding measurements---GPU overhead, real-training lead times,
large-model validation---are clearly scoped to future protocol
executions.

NRI fills a gap in the GPU systems stack that has existed since the
first transformer training run: the gap between what the training
framework produces and what a resolver can act on, filled with
structure rather than silence.

%% ─── Appendix ────────────────────────────────────────────────────────────────

\appendix

\section{Reproducibility Protocols}
\label{app:protocols}

\begin{description}[leftmargin=*,itemsep=4pt]
\item[P1 \texttt{run\_unified\_identity\_torch.py}.]
  Verifies Proposition~\ref{prop:unified} via PyTorch.
  Expected: max error $\leq 5 \times 10^{-7}$ over 10K samples.

\item[P2 \texttt{run\_overhead\_torch.py}.]
  Measures entropy overhead vs.\ attention forward pass on your
  hardware.
  Configurable device (\texttt{cpu}/\texttt{cuda}), n\_iter.
  Reports per-configuration and summary overhead percentages.
  Results should be submitted in the JSON format defined in v0.5.0.

\item[P3 \texttt{run\_early\_warning\_torch.py}.]
  Trains a small transformer in FP16 (no AMP loss scaling, high LR)
  until \texttt{NaN}.
  Reports $T_\text{alarm}$, $T_\text{NaN}$, and lead time.
  Tunable via \texttt{--lr} and \texttt{--alarm\_threshold}.

\item[P4 \texttt{run\_gpt2\_attention\_monitor.py}.]
  Loads GPT-2 small (117M), runs a calibration pass on a
  user-supplied corpus, monitors attention entropy layer-by-layer.
  Supports optional collapse injection for detection verification.
  Requires: \texttt{transformers} $\geq 4.30$.
\end{description}

\section{Library API Reference}
\label{app:api}

\begin{verbatim}
import nonans

# Training-time: single integration line
model = nonans.wrap(model, optimizer)

# Inference-time monitoring
monitor = nonans.RuntimeMonitor(sequence_length=512)
# calibration pass...
monitor.finalize_calibration()
state, conf, H = monitor.classify_head(
    attn_weights, layer=0, head=0)

# Telemetry
nonans.bus.subscribe(
    nonans.InMemorySink(capacity=1024))

# Custom resolver
class MyResolver(nonans.ResolverProtocol):
    def can_resolve(self, tensor, ctx): ...
    def resolve(self, tensor, ctx,
                budget_ms=5.0): ...
\end{verbatim}

%% ─── Citation Recommendation ─────────────────────────────────────────────────
\section{Recommended Citation}
\label{app:citation}

\textbf{Plain text (author-year):}

\begin{quote}\small
Makhebi Ahlem (2026). Numerical Runtime Intelligence (NRI): A Framework
for Real-Time Observation, Classification, and Selective Stabilization
of Numerical Instability Propagation in GPU Transformer Runtimes.
\textit{Preprint}, version 1.0.0.
DOI: \texttt{10.5281/zenodo.20573423}.
Repository: \url{https://github.com/makheahlem/nonans}.
\end{quote}

\textbf{BibTeX:}

\begin{verbatim}
@article{makhebi2026nri,
  title   = {Numerical Runtime Intelligence ({NRI}):
             A Framework for Real-Time Observation,
             Classification, and Selective Stabilization
             of Numerical Instability Propagation in
             {GPU} Transformer Runtimes},
  author  = {Makhebi, Ahlem},
  journal = {Preprint},
  year    = {2026},
  version = {1.0.0},
  doi     = {10.5281/zenodo.20573423},
  url     = {https://github.com/makheahlem/nonans}
}
\end{verbatim}

\textbf{Canonical citation wording for derivative work:}

\begin{quote}\small
``The Numerical Runtime Intelligence (NRI) framework and its
\texttt{nonans} reference implementation were introduced by
Makhebi Ahlem (2026).
The signal score, ghost/real fault taxonomy, FaultContext interface,
and three-gate monitoring architecture originated in that work.''
\end{quote}

%% ─── References ──────────────────────────────────────────────────────────────

\bibliographystyle{plainnat}
\begin{thebibliography}{19}

\bibitem[Arora et~al.(2019)]{arora2019implicit}
Arora, S., Cohen, N., Hu, W., and Luo, Y. (2019).
Implicit regularization in deep matrix factorization.
\emph{Advances in Neural Information Processing Systems (NeurIPS)}, 32:7413--7424.

\bibitem[Dong et~al.(2021)]{dong2021attention}
Dong, Y., Cordonnier, J.-B., and Loukas, A. (2021).
Attention is not all you need: pure attention loses rank exponentially with depth.
\emph{Proceedings of the 38th International Conference on Machine Learning (ICML)}, PMLR 139:2793--2803.

\bibitem[Goodfellow et~al.(2016)]{goodfellow2016deep}
Goodfellow, I., Bengio, Y., and Courville, A. (2016).
\emph{Deep Learning}.
MIT Press, Cambridge, MA. ISBN 9780262035613.

\bibitem[Halko et~al.(2011)]{halko2011finding}
Halko, N., Martinsson, P.-G., and Tropp, J.~A. (2011).
Finding structure with randomness: probabilistic algorithms for constructing approximate matrix decompositions.
\emph{SIAM Review}, 53(2):217--288.
DOI: 10.1137/090771806.

\bibitem[Jain and Wallace(2019)]{jain2019attention}
Jain, S. and Wallace, B.~C. (2019).
Attention is not explanation.
\emph{Proceedings of the 2019 Conference of the North American Chapter of the Association for Computational Linguistics: Human Language Technologies (NAACL-HLT)}, 3543--3556.
DOI: 10.18653/v1/N19-1357.

\bibitem[Kloberdanz et~al.(2022)]{kloberdanz2022deepstability}
Kloberdanz, E., Kloberdanz, K.~G., and Le, W. (2022).
DeepStability: a study of unstable numerical methods and their solutions in deep learning.
\emph{Proceedings of the 44th International Conference on Software Engineering (ICSE)}, 586--597.
DOI: 10.1145/3510003.3510095.

\bibitem[Liu et~al.(2025)]{liu2025automated}
Liu, J., Zhang, X., Xu, L., Fang, C., Gu, M., Luo, W., Chai, D., Wang, J., Zhao, Z., and Chen, Z. (2025).
Automated detection and repair of floating-point precision problems in convolutional neural network operators.
\emph{ACM Transactions on Software Engineering and Methodology}, 34(7), Article 203, 203:1--203:32.
DOI: 10.1145/3715104.

\bibitem[Manakul et~al.(2023)]{manakul2023selfcheckgpt}
Manakul, P., Liusie, A., and Gales, M.~J.~F. (2023).
SelfCheckGPT: zero-resource black-box hallucination detection for generative large language models.
\emph{Proceedings of the 2023 Conference on Empirical Methods in Natural Language Processing (EMNLP)}, 9004--9017.
DOI: 10.18653/v1/2023.emnlp-main.557.

\bibitem[Martin and Mahoney(2021)]{martin2021implicit}
Martin, C.~H. and Mahoney, M.~W. (2021).
Implicit self-regularization in deep neural networks: evidence from random matrix theory and implications for learning.
\emph{Journal of Machine Learning Research}, 22(165):1--73.

\bibitem[Micikevicius et~al.(2018)]{micikevicius2018mixed}
Micikevicius, P., Narang, S., Alben, J., Diamos, G., Elsen, E., Garcia, B., Ginsburg, B., Houston, M., Kuchaiev, O., Venkatesh, G., and Wu, H. (2018).
Mixed precision training.
\emph{International Conference on Learning Representations (ICLR)}.
URL: \url{https://openreview.net/forum?id=r1vgEWe0_}.

\bibitem[Miyato et~al.(2018)]{miyato2018spectral}
Miyato, T., Kataoka, T., Koyama, M., and Yoshida, Y. (2018).
Spectral normalization for generative adversarial networks.
\emph{International Conference on Learning Representations (ICLR)}.

\bibitem[Rodriguez et~al.(2022)]{rodriguez2022bwr}
Rodriguez, E., Alvarez-Ramirez, J., and Espinosa-Paredes, G. (2022).
A singular value decomposition entropy approach to instability analysis in BWRs.
\emph{Nuclear Engineering and Design}, 386:111559.
DOI: 10.1016/j.nucengdes.2021.111559.

\bibitem[Shannon(1948)]{shannon1948mathematical}
Shannon, C.~E. (1948).
A mathematical theory of communication.
\emph{The Bell System Technical Journal}, 27(3):379--423.
DOI: 10.1002/j.1538-7305.1948.tb01338.x.

\bibitem[Vaswani et~al.(2017)]{vaswani2017attention}
Vaswani, A., Shazeer, N., Parmar, N., Uszkoreit, J., Jones, L., Gomez, A.~N., Kaiser, {\L}., and Polosukhin, I. (2017).
Attention is all you need.
\emph{Advances in Neural Information Processing Systems (NeurIPS)}, 30:5998--6008.

\bibitem[Wiegreffe and Pinter(2019)]{wiegreffe2019attention}
Wiegreffe, S. and Pinter, Y. (2019).
Attention is not not explanation.
\emph{Proceedings of the 2019 Conference on Empirical Methods in Natural Language Processing and the 9th International Joint Conference on Natural Language Processing (EMNLP-IJCNLP)}, 11--20.
DOI: 10.18653/v1/D19-1002.

\bibitem[Yang and Dong(2024)]{yang2024neural}
Yang, X. and Dong, J.~S. (2024).
Neural Surveillance: live-update visualization of latent training dynamics.
\emph{arXiv preprint arXiv:2405.15135}.

\bibitem[Yunis et~al.(2024)]{yunis2024spectral}
Yunis, D., Patel, K.~K., Wheeler, S., Savarese, P., Vardi, G., Livescu, K., Maire, M., and Walter, M.~R. (2024).
Approaching deep learning through the spectral dynamics of weights.
\emph{Advances in Neural Information Processing Systems (NeurIPS)}, 37:42011--42035.

\bibitem[Zhai et~al.(2023)]{zhai2023sigma}
Zhai, S., Likhomanenko, T., Littwin, E., Busbridge, D., Ramapuram, J., Zhang, Y., Gu, J., and Susskind, J.~M. (2023).
Stabilizing transformer training by preventing attention entropy collapse.
\emph{Proceedings of the 40th International Conference on Machine Learning (ICML)}, PMLR 202:40770--40803.

\bibitem[Zhang et~al.(2022)]{zhang2022wse}
Zhang, M., Kong, P., Hou, A., Xia, A., Tuo, W., and Lv, Y. (2022).
Identification strategy design with the solution of wavelet singular spectral entropy algorithm for the aerodynamic system instability.
\emph{Aerospace}, 9(6):320.
DOI: 10.3390/aerospace9060320.

\end{thebibliography}

\end{document}
